\section{Trace-Length Optimization}
\label{app:trace-length-optimization}

This appendix evaluates the trace quantities whose exact interfaces appear in
Section~\ref{sec:trace-length}.  The center and V triangle coordinates are
used only here; the body retains the geometric hypotheses and final bounds.

\subsection{Conditional skeleton caps}

\begin{lemma}[Center-skeleton calculation]
\label{lem:app-center-skeleton-cap-calculation}
Under the hypotheses of Lemma~\ref{lem:center-skeleton-cap}, one has
$L_S(T_C)<3/2$.
\end{lemma}

\begin{proof}
Normalize the unique midpoint as $M_0$ and a positive boundary trace to
$e_{0,1}$.  Use the signed parameters and six radial exits of
\zcref{prop:signed-center-normal-form}. The positive trace condition is
\begin{equation}
 \alpha+W\delta<P.
 \label{eq:center-skeleton-domain}
\end{equation}
Summing the boundary and radial contributions gives
\[
 L_S(T_C)\le K_R+\mathcal E(\alpha,\delta),
 \qquad K_R:=W+E-\frac\eta R,
\]
where
\[
 \mathcal E(\alpha,\delta)
 :=-\frac\alpha R+\frac\alpha W+\delta
 +\min\left\{\frac\alpha R,\frac\delta W\right\}
 +\frac{[P-(R\alpha+\delta)]_+}{W}.
\]
Splitting according to the minimum and positive part gives
\begin{equation}
 \mathcal E(\alpha,\delta)\le\frac PW.
 \label{eq:center-skeleton-E-bound}
\end{equation}
If $\alpha/R\le\delta/W$, the expression without the positive part is
$\alpha/W+\delta$; when that part is active, the excess over $P/W$ is
$\alpha-R\delta/W\le0$.  The reflected order gives
$\delta-W\alpha/R\le0$.  If the positive part vanishes,
\eqref{eq:center-skeleton-domain} gives the same bound strictly.

Put $A:=1+R-R^2$.  Then
\[
 \frac32-K_R-\frac PW=\frac{1+A-2EA}{2RW}.
\]
Both sides of $1+A>2EA$ are positive, and
\[
 (1+A)^2-4A^2E^2=R^2W^2(5+4R-4R^2)>0.
\]
Therefore $K_R+P/W<3/2$, as required.
\end{proof}


For the remaining cap, reflect the local directions if necessary so that
$0\le A_i\le B_i\le1$.  The actual reach triple of a supercritical V
triangle lies in the selected cell
\begin{equation}
\begin{aligned}
 A_i^2+A_iB_i+B_i^2&\le1,\\
 C_i&\le\frac12,\\
 (A_i^2-1)C_i^2+(2A_iB_i^2+B_i)C_i+(B_i^4-B_i^2)&\le0,
\end{aligned}
\label{eq:supercritical-admissible-cell}
\end{equation}
where $C_i$ belongs to the connected component containing $0$ in the
corresponding one-dimensional fiber.  This is Cell $S$ of
Proposition~\ref{prop:exact-admissible-set}.  The selector is essential: the
quadratic is negative at $C_i=0$, positive at $C_i=1/2$, and opens downward.

\begin{lemma}[Supercritical skeleton calculation]
\label{lem:app-supercritical-skeleton-calculation}
Every V triangle with $A_i+B_i>1$ has $L_S(T_i)<3/2$.
\end{lemma}

\begin{proof}
The corner normal form, Proposition~\ref{prop:vd-corner-normal-form}, makes
$\Vdone$ and $\Vdtwo$ triangles nonsupercritical.  Lemma
\ref{lem:t3-nonsupercritical} gives the same conclusion for $\Tlike$
triangles.  Hence a supercritical triangle is $\Vdzero$ and has no positive
adjacent-radial support.  Reflect so that $A_i\le B_i$.
Then
\[
 L_S(T_i)=A_i+B_i+C_i.
\]

Set
\[
 s:=A_i+B_i,\qquad \delta:=B_i-A_i,
 \qquad \gamma_0:=\frac32-s.
\]
The diameter constraint gives
\[
 1<s\le\frac2{\sqrt3},
 \qquad0\le\delta\le\sqrt{4-3s^2}.
\]
Let $P(\gamma)$ be the quadratic in
\eqref{eq:supercritical-admissible-cell}.  Substitution and multiplication by
$16$ yield
\[
\begin{aligned}
16P(\gamma_0)=Q(s,\delta)
&=\delta^4+8\delta^3s-6\delta^3
 +14\delta^2s^2-18\delta^2s+5\delta^2\\
&\quad-8\delta s^3+30\delta s^2-34\delta s+12\delta\\
&\quad+s^4-6s^3-19s^2+60s-36.
\end{aligned}
\]
Its derivative factors as
\[
 \frac{\partial Q}{\partial\delta}
 =2(2\delta+4s-3)
 \bigl(\delta^2+\delta(4s-3)+(s-1)(2-s)\bigr)\ge0.
\]
Therefore
\[
 Q(s,\delta)\ge Q(s,0)
 =(s-1)(s^3-5s^2-24s+36)>0.
\]
The cubic decreases on $[1,2/\sqrt3]$ and at the right endpoint equals
$88/3-136\sqrt3/9>0$.  Hence $P(\gamma_0)>0$.  The selected smaller-root
component in \eqref{eq:supercritical-admissible-cell} forces
$C_i<\gamma_0$, so $A_i+B_i+C_i<3/2$.
\end{proof}


\subsection{Vd1 and Vd2 corner optimization}

Use temporary coordinates
\[
 X=V_0+x(V_5-V_0)+y(V_1-V_0),
\]
so $V_0=(0,0)$, $V_5=(1,0)$, $V_1=(0,1)$, $O=(1,1)$, and
$\lVert(x,y)\rVert^2=x^2+y^2-xy$.

\begin{proposition}[Vd1/Vd2 corner normal form]
\label{prop:vd-corner-normal-form}
Let $T$ be the closure of an original $\Vdone$ or $\Vdtwo$ triangle at
$V_0$, and let $A_0,B_0$ be its exact reaches toward $V_5,V_1$.  There is a
unique $t>0$ such that, with $d=\sqrt{t^2+t+1}$,
\begin{equation}
 T=\left\{(x,y):
 \begin{aligned}
 x-(t+1)y&\le A_0,\\
 ty-(t+1)x&\le tB_0,\\
 tx+y&\le d-A_0-tB_0
 \end{aligned}\right\}.
 \label{eq:vd-normal-form}
\end{equation}
The raw traces, parametrized from the indicated outer vertex, are
\begin{equation}
\begin{array}{c|c}
r_0&0\le s\le\dfrac{d-A_0-tB_0}{t+1}\\[6pt]
r_1&\dfrac{t(1-B_0)}{t+1}\le s\le\dfrac{d-A_0-tB_0-1}{t}\\[8pt]
r_5&\dfrac{1-A_0}{t+1}\le s\le d-A_0-tB_0-t.
\end{array}
\label{eq:vd-ray-intervals}
\end{equation}
Intersecting with $[0,1]$ gives the actual traces, and
\begin{equation}
 A_0+B_0<\frac12.
 \label{eq:vd-half-cap}
\end{equation}
\end{proposition}

\begin{proof}
The triangle has one vertex $W$ outside $H$ and two vertices $U,Z$ in $H$.
Because $V_0$ is an interior convex combination, both coordinates of $W$ are
negative.  The axis endpoints $(A_0,0),(0,B_0)$ lie on the two sides through
$W$.  Put $p=U-W$ and $q=Z-W$.  Their order selects
\[
 q=(p_x-p_y,p_x),\qquad p_x>p_y>0.
\]
Writing $t=(p_x-p_y)/p_y$ gives uniquely
\[
 p=\frac{(t+1,1)}d,\qquad q=\frac{(t,t+1)}d,
\]
which yields \eqref{eq:vd-normal-form}; substitution of
$(s,s),(s,1),(1,s)$ yields \eqref{eq:vd-ray-intervals}.

If $r_1$ has positive trace, clearing its endpoint inequality gives
\[
 (t+1)A_0+tB_0<(t+1)(d-1)-t^2.
\]
Since the left side is at least $t(A_0+B_0)$,
\[
 A_0+B_0<\frac{(t+1)(d-1)-t^2}{t}<\frac12.
\]
The last inequality follows from
\[
 \left(t^2+\frac{3t}{2}+1\right)^2
 -(t+1)^2(t^2+t+1)=\frac{t^2}{4}>0.
\]
Reflection treats support on $r_5$.
\end{proof}

\begin{lemma}[Vd2 adjacent-midpoint cap]
\label{lem:vd2-neighbor-midpoint-cap}
If an original open V triangle has $\Vdtwo$ closure $T_i$ and
\[
 \{M_{i-1},M_{i+1}\}\cap T_i\ne\varnothing,
\]
then its exact boundary reaches satisfy
\[
 A_i+B_i<\frac13.
\]
\end{lemma}

\begin{proof}
In Proposition~\ref{prop:vd-corner-normal-form}, reflection sends
\[
 (t,A_0,B_0,M_1,M_5)
 \longmapsto(1/t,B_0,A_0,M_5,M_1).
\]
Thus assume $t\ge1$ and use $A_0+B_0\le A_0+tB_0$.

If $M_5=(1,1/2)$ lies in $T$, the third half-plane gives
\[
 A_0+B_0\le A_0+tB_0
 \le d-t-\frac12\le\sqrt3-\frac32<\frac13.
\]
If $M_1=(1/2,1)$ lies in $T$, the second and third half-planes give
\[
 B_0\ge\frac{t-1}{2t},
 \qquad A_0+B_0\le S(t):=d-t-\frac1{2t}.
\]
Positive-length contact with $r_5$ gives
\[
 A_0+B_0<R(t):=\frac{2(t+1)d-3t^2-t-2}{2t}.
\]
Direct differentiation shows $S$ increasing and $R$ decreasing on
$[1,\infty)$.  For $S'$, the sign reduces to
\[
 (2t+1)^2t^4-(t^2+t+1)(2t^2-1)^2
 =(t+1)(t^3+3t^2-1)>0,
\]
while $R'<0$ follows from $d>t+1/2$.  Splitting at $t=4/3$ gives
\[
\begin{aligned}
1\le t\le\frac43&:\quad
A_0+B_0\le\frac{8\sqrt{37}-41}{24}<\frac13,\\
t\ge\frac43&:\quad
A_0+B_0<\frac{7\sqrt{37}-39}{12}<\frac13.
\end{aligned}
\]
The final comparisons are $8\sqrt{37}<49$ and $7\sqrt{37}<43$.
\end{proof}
